% Standalone Section VI subsection on the cost of quantization (M3 + M6).
% Supersedes the Section VI half of sec4f_sec6_quant_draft.tex; that file's
% Section IV-F paragraph (what is actually quantized) still stands.
% Sources: results/tables/tab_m3_quant.{csv,tex}, m3_quant_per_seed.csv,
%   results/diagnostics/fc1_ratio_repair.csv,
%   results/diagnostics/leakage_audit/README.md
% All eleven datasets are on repaired corpora. Generated 2026-09-21.
% Labels to verify: sec:results, subsec:audit, tab:m3:quant, subsec:iot (VI-A).

\subsection{What Quantization Costs, and Why It Is Not Uniform}
\label{subsec:quantcost}

Reporting a quantized artifact obliges us to say what the projection costs.
Table~\ref{tab:m3:quant} gives fp32 and int8-grid test macro-$F_1$ for all four
models on every benchmark as 5-seed means, with the mean and worst-seed
difference. For the two tree-guided models the floating-point side comes from
re-running training with the same seeds and capturing the weights immediately
before the projection; every such re-run reproduces its int8 number exactly, so
the two columns differ only in the projection. Vanilla SNN and TASNN are
reported in floating point throughout the paper, so their int8 columns are
counterfactual. For those two we also project every weight matrix rather than
\texttt{Linear} layers alone, because Vanilla SNN keeps 98\% of its weights in a
recurrent cell that a \texttt{Linear}-only projection never touches; for the
tree-guided students, whose weights are all linear, the two variants coincide.

\textbf{The projection is close to free for the spiking student and not for the
others.} Tree-to-Spike loses $0.51$ macro-$F_1$ points on Edge-IIoTset and
$1.58$ on CSE-CIC-IDS2018, and averages $+0.06$ across the nine remaining
benchmarks; its worst single seed anywhere is $-4.38$. Its dense twin loses
$5.97$ on Edge-IIoTset, with one seed at $-13.20$. The counterfactual for
Vanilla SNN is larger still: $-14.91$ on Edge-IIoTset and $-20.18$ at its worst
seed, so the baseline we did \emph{not} deploy at int8 is the model the
projection would have hurt most. Post-training per-tensor quantization is
therefore not automatically safe, and reporting a quantized artifact without
measuring it is not a neutral omission.

\textbf{The mechanism is a single outlier weight.} The scale is
$s = \max|W|/127$ per tensor, so one large weight fixes the resolution of every
other weight in that tensor. Where the tree-guided students' per-layer
$\max|W|/p_{99}|W|$ ratios stay between 1.1 and 3.6, the dense twin's first
layer -- the one reading the encoder -- reaches 8.5--9.8 on three of five
CSE-CIC-IDS2018 seeds and 6.1 on the worst Edge-IIoTset seed, where 42\% of its
first-layer weights round to zero. Accuracy tracks that ratio: the seeds with
ratios near 2.5--3.0 lose 0.2 points or less, the seeds near 9 lose 1.5--3.6.
The spiking student does not accumulate comparable outliers on the same encoder
input, which is what makes its deployed numbers insensitive to the projection.

\textbf{Two causes, one amplifying the other.} Before the dataset repair of
Section~\ref{subsec:audit}, this cell was far worse: DNN-Tree lost $11.46$ points
on CSE-CIC-IDS2018, with a worst seed of $-29.27$ and a first-layer ratio of
15.9. The leaked \texttt{Src~Port} feature was the teacher's most important
input, and the outlier weight sat on it; removing the artifact removed the
amplification but not the fragility, which now shows up on Edge-IIoTset instead.
The dense twin has a per-tensor quantization fragility of its own, and a dataset
artifact that concentrates importance in one feature makes it an order of
magnitude worse. Neither explanation alone fits the measurements.

% Section VI-A (per-dataset results), where the Edge-IIoTset ordering is given:
\paragraph{Precision changes the Edge-IIoTset ordering.}
On Edge-IIoTset the deployed (int8) ordering places Tree-to-Spike at 83.95
macro-$F_1$ and DNN-Tree at 80.83, a 3.1-point lead; in floating point the same
runs give 84.46 and 86.81, a 2.4-point lead the other way. The reversal is
entirely the quantization loss of Table~\ref{tab:m3:quant}. We report the
deployed artifacts throughout, because the claim of this paper concerns what is
shipped to a constrained device, and on that basis the dense twin's fragility is
a genuine deployment cost rather than a measurement artifact. A reader comparing
the two \emph{architectures} at equal precision should use the fp32 columns,
where the dense twin is ahead on this dataset.
